\section{Architecture--Demand Lattice}
\label{sec:lattice}

\begin{figure*}[t]
  \centering
  \paperfigure{figure2_control_lattice}{\textwidth}
  \caption{\textbf{Matched adjacent gains in static and repeated controlled
  probes.} Panel A shows static E12. Panel B shows E18b's registered mean
  affected-MSE gain over the update-by-gap grid. Bars are frozen seed means
  and dots are paired seed-level gains. E18b passed its relative cell-mean
  simpler-demand and retention guards, and the registered stress-cell
  direction was positive in 5/5 paired seeds. E18 supplies oracle addresses,
  candidates, demand descriptors, and a model-visible verified-event bit. The
  result does not establish gain in every cell, stress SESOI, absolute
  preservation, minimal-controller sufficiency, semantic demand inference, or
  learned relevance.}
  \label{fig:lattice}
\end{figure*}

\subsection{Static controlled lattice}

Table~\ref{tab:lattice} maps each added freedom to the demand constructed to
require it.  The frozen E12 registered-grid gains are
$\CATENAResult{E12_MAGNITUDE_GAIN}$,
$\CATENAResult{E12_VALUE_GAIN}$,
$\CATENAResult{E12_ADDRESS_GAIN}$, and
$\CATENAResult{E12_STATE_GAIN}$ for the magnitude, value, address, and
state-conditioned transitions, respectively.  The inferential point is the
matched adjacent gain together with the registered guardrails, not raw
controller size.

\begin{table*}[t]
  \centering
  \small
  \begin{tabular}{llll}
    \toprule
    Added freedom & Matched demand & Adjacent contrast & Guardrail scope\\
    \midrule
    Erase/write magnitude & asymmetric same address
      & tied $\rightarrow$ dual & symmetric demands, retention\\
    Value channels & partial value subspace
      & dual $\rightarrow$ diagonal & scalar-demand non-inferiority\\
    Address decoupling & different erase/write addresses
      & diagonal $\rightarrow$ separate & prior-demand non-inferiority\\
    State conditioning & state-dependent target
      & separate $\rightarrow$ state-aware & prior-demand non-inferiority\\
    \bottomrule
  \end{tabular}
  \caption{The registered architecture--demand map. ``Selective'' refers to a
  matched adjacent gain plus the corresponding simpler-demand and retention
  guardrails; it does not assert zero benefit on every other harder demand.}
  \label{tab:lattice}
\end{table*}

\subsection{Repeated-sequence lattice}

E18 repeats the nested comparison over update-count and distractor-gap cells
while supplying oracle erase/write addresses, oracle candidates, explicit
demand descriptors, and a visible verified-event bit.  Its registered-grid
mean adjacent gains are
$\CATENAResult{E18_MAGNITUDE_GAIN}$,
$\CATENAResult{E18_VALUE_GAIN}$,
$\CATENAResult{E18_ADDRESS_GAIN}$, and
$\CATENAResult{E18_STATE_GAIN}$.  All four registered adjacent conjunctions
passed.  Simpler-demand and retention checks are maximum relative cell-mean
non-inferiority guards, not absolute-accuracy claims.  At the registered
8-update, 2,048-distractor stress cell, the direction was positive in 5/5
paired seeds; there was no separate stress SESOI.
